\section{Mechanisms}
\label{sec:mechanisms}

The results establish that first-time lucky loser entry causes gains in ranking points and tournament access. This section decomposes the effect into two channels, access and ability, and examines match-level performance conditional on opponent strength.

\subsection{Access versus Ability}
\label{sec:mechanisms_access}

The career effects of LL entry operate primarily through tournament access rather than improved competitive skill. On the ATP tour, ranking points rise significantly through 26 weeks while Elo ratings are null at every horizon. Because Elo adjusts for opponent quality, the null Elo result indicates that treated players do not perform better or worse per match once opponent strength is held constant. The ranking-point gains come entirely from additional tournament participation: LL entry provides a main-draw appearance that generates ranking points, which improve the player's position in the rolling 52-week ranking, which crosses direct-entry thresholds for subsequent lower-tier events. The evidence aligns with this access cascade: treated players enter more main draws and play more matches at 250+ events over 26 weeks, and their quality-adjusted match performance is unchanged.

The WTA tour presents a nuance. Elo gains at 8 weeks ($+15.0$, $p = 0.082$) and 52 weeks ($+36.8$, $p = 0.055$) are consistent with WTA LL recipients experiencing modest ability improvements. The same pattern appears in the non-GS WTA results, where Elo gains reach significance at 12 weeks ($+10.5$, $p = 0.011$). Greater competitive parity on the women's tour would give marginal players more room to accumulate ability gains through main-draw experience against stronger opponents. The match-level tournament performance model (Section~\ref{sec:mechanisms_performance}) does not detect a significant WTA effect on win probability conditional on opponent strength, however, so the ability interpretation carries only limited weight. The tension between Elo gains in the stacked model and the null match-level win probability estimate ($\hat{\delta} = +0.146$, $p = 0.29$; player-clustered) can plausibly reflect Elo's sensitivity to the schedule of opponents faced. If treated WTA players enter more tournaments with weaker fields (lower-tier events where ranking gains provide access), Elo can improve even without a genuine ability change, because the Elo system rewards wins against any opponent. The match-level model conditions on actual opponent strength and provides a stricter test of ability improvement.

The access cascade is self-limiting. As the 52-week rolling window progresses, the points from the original LL event constitute a shrinking fraction of the player's total. Unless downstream tournaments generate enough additional points to replace the original ones, the ranking converges back toward the level determined by sustained match performance. The dynamic results align with this: ATP ranking-point gains peak at 12 weeks and fade toward zero at 52 weeks.

\subsection{Match-Level Performance}
\label{sec:mechanisms_performance}

To test directly whether LL entry affects competitive performance itself, I estimate a match-level logit of win probability conditional on opponent strength and tournament context (see Appendix~\ref{app:match_model}). At Grand Slams, the lottery provides exogenous variation; at non-GS events, a control function corrects for endogenous selection.

First-time ATP GS LL recipients have $\hat{\delta} = -0.141$ (SE $= 0.103$, $p = 0.17$; player-clustered), a point estimate consistent with overseeding when novice entrants face opponents above their ability level, though not significant at conventional levels. WTA GS recipients have $\hat{\delta} = +0.146$ (SE $= 0.138$, $p = 0.29$), also not statistically distinguishable from zero. Non-GS ATP has $\hat{\delta} = +0.609$ ($p < 0.001$), reflecting the control-function correction in a setting where treated players are positively selected on ability. The GS point estimates line up with the direction the overseeding channel predicts (negative for ATP, positive for WTA), but the small first-time GS samples ($N = 120$ ATP, $82$ WTA) do not deliver enough power for individually significant tests. Full results including dose and horizon interactions are reported in Appendix Table~\ref{tab:delta_all}.

The event study figures (Figures~\ref{fig:event_study_atp} and~\ref{fig:event_study_wta}) provide a visual summary: ranking-point trajectories diverge sharply after the LL event while pre-event trends are parallel, so the treatment effect is concentrated in post-event access gains rather than pre-existing ability differences.

\subsection{Economic Magnitude: Prize Money}
\label{sec:mechanisms_prize}

Ranking points are an imperfect proxy for the economic stakes. To ground the results in dollar terms, I compute the implied prize-money differential using publicly available tournament prize tables. At Grand Slams, a first-time LL recipient earns \$100{,}000--\$400{,}000 USD depending on how far they advance in the main draw, versus a \$40{,}000 qualifying-round stipend for controls. Combined with the estimated access gains over the 26 weeks following the focal event (evaluated at an average forward-looking first-round prize of \$11{,}000 per additional main draw), the total prize-money effect of a first-time LL entry is approximately \$101{,}000 on the ATP GS sample and \$94{,}000 on the WTA GS sample (Appendix Table~\ref{tab:prize}). For non-GS first-time LL recipients, the total dollar effect is smaller but economically meaningful: \$28{,}000 on the ATP tour and \$16{,}000 on the WTA tour, driven more heavily by the dynamic access cascade than by the immediate event prize. These magnitudes are substantial against qualifying-level players' annual earnings, which typically fall in the \$40{,}000--\$80{,}000 range. Under a $\pm 20\%$ sensitivity band on the per-entry prize assumption (\$8{,}800--\$13{,}200 instead of \$11{,}000), the ATP GS total effect ranges from approximately \$99{,}000 to \$103{,}000 and the non-GS ATP effect from approximately \$24{,}000 to \$31{,}000; the qualitative conclusion (dollar effects comparable to a qualifier's annual earnings) is insensitive to reasonable variation in the per-entry assumption. The first break therefore carries real financial consequences even as its ranking-point effects fade within a year.

\subsection{Reconciling with the Existing Literature}
\label{sec:mechanisms_reconciling}

\citet{Maity2025_early_career_sports} report positive correlations between lucky loser entry and career outcomes. The match-level tournament performance model, which exploits the lottery for identification, does not detect an improvement in match-level win probability on the ATP tour and reports only a marginal positive effect on the WTA tour. The causal estimates here indicate that the correlational finding in \citet{Maity2025_early_career_sports} likely reflects the additional tournament participation that LL entry generates rather than an improvement in competitive performance itself.

The initial conditions literature reports persistent effects of early-career shocks \citep{Oyer2006_initial_placement, Oreopoulos2012_graduating_recession}. The present results align on the access dimension: LL entry causes persistent changes in the number and type of tournaments a player enters; the ability dimension is largely null. Market transparency is the plausible moderator: in professional tennis, where performance is measured match by match and rankings update weekly, ability reasserts itself within a year. The access channel alone accounts for the ranking-point gains and tournament-entry effects documented in Section~\ref{sec:results}.
